\chapter*{Chapter Summaries}

% One short paragraph per chapter, kept in sync with the actual content.

\paragraph{Chapter 1 — Introduction.} Motivates the use of small language
models (SLMs) as web agents, states the problem of automatically generating
high-quality training data from a website without human annotation, and lists
the objectives, contributions, and scope of the thesis.

\paragraph{Chapter 2 — Background and Related Work.} Reviews the web-agent
task formulation and the WebArena benchmark, then surveys LLM-based web
agents, skill and workflow discovery, automatic training-data generation, and
small language models as agents, positioning this thesis against each line of
work.

\paragraph{Chapter 3 — Proposed Method.} Presents HSGE, a graph-based
exploration pipeline: a two-level graph representation of a website, graph-guided
exploration with LLM-proposed low-level tasks, skill discovery and
verification, semantic task generation by chaining planning operators, data
curation, and LoRA fine-tuning of an SLM.

\paragraph{Chapter 4 — Experiments.} Describes the experimental setup on
self-hosted WebArena sites, reports data-collection and fine-tuning results,
and evaluates the fine-tuned SLM against its base model on the official
WebArena task suites of four sites: Shopping, Shopping Admin, GitLab, and
Reddit.

\paragraph{Chapter 5 — Discussion.} Interprets the results and distills the
engineering experience gained while running LLM agents on live websites into
reusable lessons and best practices, together with threats to validity.

\paragraph{Chapter 6 — Conclusion.} Summarizes the products and contributions,
compares the achieved results against the initial objectives, and outlines
future research directions.
